\documentclass[11pt]{article}

\usepackage[utf8]{inputenc}
\usepackage[T1]{fontenc}
\usepackage{lmodern}
\usepackage[margin=1in]{geometry}
\usepackage{amsmath,amssymb,mathtools,bm}
\usepackage{booktabs}
\usepackage{enumitem}
\usepackage{microtype}
\usepackage{hyperref}
\usepackage{xcolor}
\usepackage{listings}
\usepackage{longtable}
\usepackage{array}

\hypersetup{
  colorlinks=true,
  linkcolor=blue,
  citecolor=blue,
  urlcolor=blue
}

\lstset{
  basicstyle=\ttfamily\small,
  columns=fullflexible,
  frame=single,
  breaklines=true,
  showstringspaces=false,
  keywordstyle=\color{blue!70!black},
  commentstyle=\color{green!40!black}
}

\newcommand{\R}{\mathbb{R}}
\newcommand{\dd}{\mathrm{d}}
\newcommand{\T}{\mathsf{T}}
\newcommand{\grad}{\nabla}
\newcommand{\Hess}{\nabla^2}
\newcommand{\Id}{\mathbf{I}}
\newcommand{\zero}{\mathbf{0}}
\newcommand{\x}{\mathbf{x}}
\newcommand{\vct}{\mathbf{v}}
\newcommand{\aV}{\mathbf{a}}
\newcommand{\mV}{\mathbf{m}}
\newcommand{\gV}{\mathbf{g}}
\newcommand{\Q}{\mathbf{Q}}
\newcommand{\Mmat}{\mathbf{R}}
\newcommand{\Kmat}{\mathbf{K}}
\newcommand{\Jmat}{\mathbf{J}}
\newcommand{\Pmat}{\mathbf{P}}
\newcommand{\atanTwo}{\operatorname{atan2}}
\newcommand{\diag}{\operatorname{diag}}
\newcommand{\supp}{\operatorname{supp}}

\title{Exact Discrete Bending Hessian for the Mem3DG Helfrich Energy\\
\large Mathematical derivation, geometric primitives, sparse assembly, and implementation checks}
\author{Implementation-oriented technical note}
\date{\today}

\begin{document}
\maketitle
\tableofcontents
\newpage

\section{Purpose and scope}

This document derives the exact Cartesian Hessian of the discrete spontaneous-curvature bending energy used by Mem3DG.  The intended reader is implementing one of the following:
\begin{enumerate}[label=(\alph*)]
  \item an assembled sparse Newton matrix,
  \item a matrix-free Hessian--vector product for Newton--Krylov methods, or
  \item a reference implementation based on second-order automatic differentiation.
\end{enumerate}

The central result is deliberately written in terms of four geometric derivatives for each vertex $i$:
\begin{equation}
  \aV_i = \grad A_i,
  \qquad
  \mV_i = \grad M_i,
  \qquad
  \Q_i = \Hess A_i,
  \qquad
  \Mmat_i = \Hess M_i,
\end{equation}
where $A_i$ is the barycentric dual area and $M_i=A_iH_i$ is the integrated mean curvature.  Once these four quantities are available, the bending Hessian is a short closed-form expression.  This separation is useful because it isolates the energy calculus from the more delicate geometric calculus of triangle areas and signed dihedral angles.

The derivation assumes:
\begin{itemize}
  \item a fixed triangulation and fixed topology during differentiation;
  \item nondegenerate triangles and nonzero edge lengths;
  \item a consistently oriented manifold triangle mesh;
  \item bending modulus $\kappa_i$ and spontaneous curvature $c_i$ treated as fixed vertex data during the shape derivative; and
  \item the Gaussian-curvature term omitted for a closed fixed-topology mesh with constant Gaussian modulus, because it is topologically constant.
\end{itemize}

Boundary terms, variable Gaussian modulus, geometry-dependent material fields, remeshing, and topology changes require additional terms and are discussed separately.

\section{Important normalization and sign conventions}

\subsection{Paper convention}

The Mem3DG paper uses the discrete bending energy
\begin{equation}
  E_b(\x)
  =
  \sum_{i=1}^{N}
  \kappa_i A_i(\x)\bigl(H_i(\x)-c_i\bigr)^2,
  \label{eq:paper-energy}
\end{equation}
where $c_i$ denotes the spontaneous curvature $H_i^*$ in the paper.

The paper defines pointwise mean curvature from integrated mean curvature by
\begin{equation}
  H_i = \frac{M_i}{A_i},
  \qquad
  M_i = A_iH_i
      = \frac{1}{4}\sum_{e\ni i}\ell_e\phi_e.
  \label{eq:Hi-Mi}
\end{equation}
Here $\ell_e$ is edge length and $\phi_e$ is the signed dihedral angle.  The factor $1/4$ is the combination of the edge integrated mean-curvature factor $\ell_e\phi_e/2$ and the equal split of an edge quantity between its two endpoint vertices.

\subsection{Current Mem3DG source normalization}

At the time this note was prepared, the current Mem3DG source computes the spontaneous-curvature energy in the form
\begin{equation}
  E_b^{\mathrm{code}}
  =
  2\sum_i K_{b,i}A_i(H_i-H_{0,i})^2.
  \label{eq:code-energy}
\end{equation}
Therefore all formulas below match the current source by setting
\begin{equation}
  \boxed{\kappa_i = 2K_{b,i},\qquad c_i=H_{0,i}.}
  \label{eq:normalization-map}
\end{equation}
The source writes an absolute value before squaring the curvature mismatch.  Algebraically,
$\lvert H_i-H_{0,i}\rvert^2=(H_i-H_{0,i})^2$, so the absolute value should be removed before differentiation.

This factor of two is consistent with the common continuum convention
\begin{equation}
  E=\frac{K_b}{2}\int(2H-C_0)^2\,\dd A
   =2K_b\int\left(H-\frac{C_0}{2}\right)^2\,\dd A.
\end{equation}

\subsection{Energy gradient, physical force, and force Jacobian}

This note uses
\begin{equation}
  \mathbf{g}_b = \grad_{\x}E_b,
  \qquad
  \Kmat_b = \Hess_{\x}E_b.
\end{equation}
The physical conservative force and its Jacobian are
\begin{equation}
  \boxed{
  \mathbf{f}_b=-\mathbf{g}_b,
  \qquad
  \Jmat_b^{\mathrm{force}}
  =\frac{\partial\mathbf{f}_b}{\partial\x}
  =-\Kmat_b.}
  \label{eq:force-sign}
\end{equation}
If a solver uses the stationarity residual $\mathbf{r}=\grad E$ rather than the physical force, its Newton matrix is $+\Kmat_b$.

\section{Global and local coordinates}

Let the mesh have $N$ vertices with positions $\mathbf{x}_i\in\R^3$.  Stack all Cartesian degrees of freedom as
\begin{equation}
  \x=
  \begin{bmatrix}
    \mathbf{x}_1\\ \vdots\\ \mathbf{x}_N
  \end{bmatrix}
  \in\R^{3N}.
\end{equation}

For a vertex $i$, define its closed one-ring vertex set
\begin{equation}
  S_i=\{i\}\cup N(i).
\end{equation}
Both $A_i$ and $M_i$ depend only on positions in $S_i$ for a manifold triangle fan.  Let
\begin{equation}
  \x_i^{\mathrm{loc}}=\Pmat_i\x\in\R^{3|S_i|}
\end{equation}
be the local coordinate vector, where $\Pmat_i$ is a Boolean extraction matrix.  A local Hessian is scattered globally as
\begin{equation}
  \Kmat_b
  =
  \sum_i \Pmat_i^\T\Kmat_i^{\mathrm{loc}}\Pmat_i.
  \label{eq:sparse-assembly}
\end{equation}
Although each energy contribution has one-ring support, two vertices in the same one-ring of a center vertex can be two graph edges apart.  Consequently, the assembled Hessian generally has a two-ring coupling pattern.

\section{Discrete geometric quantities}

\subsection{Barycentric dual area}

For an oriented triangle $f=(p,q,r)$, let $A_f$ be its unsigned area.  Mem3DG uses the barycentric dual area
\begin{equation}
  \boxed{
  A_i=\frac{1}{3}\sum_{f\ni i}A_f.}
  \label{eq:dual-area}
\end{equation}
This quantity is positive on a nondegenerate mesh and has dimensions of length squared.

\subsection{Integrated and pointwise mean curvature}

For an interior edge $e=(p,q)$ with length $\ell_e$ and signed dihedral angle $\phi_e$, define the edge integrated mean curvature
\begin{equation}
  M_e^{\mathrm{edge}}=\frac{1}{2}\ell_e\phi_e.
\end{equation}
Splitting equally between the two edge endpoints gives the vertex integrated mean curvature
\begin{equation}
  \boxed{
  M_i=\frac{1}{4}\sum_{e\ni i}\ell_e\phi_e.}
  \label{eq:integrated-mean-curvature}
\end{equation}
The corresponding pointwise mean curvature is
\begin{equation}
  \boxed{H_i=\frac{M_i}{A_i}.}
\end{equation}
Thus $M_i$ has dimensions of length and $H_i$ has dimensions of inverse length.

In Geometry Central and the current Mem3DG source, the quantity named \texttt{vertexMeanCurvatures} is the integrated quantity $M_i$.  The pointwise quantity is obtained by dividing by \texttt{vertexDualAreas}.

\subsection{Orientation and signed dihedral angle}

For an interior edge shared by consistently oriented faces
\begin{equation}
  f_+=(a,b,c),
  \qquad
  f_-=(b,a,d),
\end{equation}
define
\begin{align}
  \mathbf{d}_e &= \mathbf{x}_b-\mathbf{x}_a,
  &\ell_e&=\|\mathbf{d}_e\|,
  &\mathbf{t}_e&=\frac{\mathbf{d}_e}{\ell_e},\\
  \mathbf{c}_+&=(\mathbf{x}_b-\mathbf{x}_a)\times(\mathbf{x}_c-\mathbf{x}_a),
  &\mathbf{n}_+&=\frac{\mathbf{c}_+}{\|\mathbf{c}_+\|},\\
  \mathbf{c}_-&=(\mathbf{x}_a-\mathbf{x}_b)\times(\mathbf{x}_d-\mathbf{x}_b),
  &\mathbf{n}_-&=\frac{\mathbf{c}_-}{\|\mathbf{c}_-\|}.
\end{align}
A robust signed angle is
\begin{equation}
  \boxed{
  \phi_e=\atanTwo(\sigma_e,\chi_e),
  \qquad
  \sigma_e=\mathbf{t}_e\cdot(\mathbf{n}_+\times\mathbf{n}_-),
  \qquad
  \chi_e=\mathbf{n}_+\cdot\mathbf{n}_-.}
  \label{eq:signed-dihedral}
\end{equation}
For an exactly flat pair of consistently oriented faces, $\sigma_e=0$, $\chi_e=1$, and $\phi_e=0$.  A different global sign convention is acceptable, but the energy, force, spontaneous curvature, and all derivative formulas must use the same convention.

Avoid allowing $\phi_e$ to jump across the $\pm\pi$ branch cut during a Newton step.  Near a fold crossing the branch cut, use step control or an angle-unwrapping strategy.

\section{Vertex energy in a quotient-free form}

Define the curvature mismatch
\begin{equation}
  d_i=H_i-c_i.
\end{equation}
The vertex energy is
\begin{equation}
  E_i=\kappa_iA_i d_i^2.
\end{equation}
Using $H_i=M_i/A_i$, this can be written as
\begin{equation}
  \boxed{
  E_i
  =\kappa_i\frac{(M_i-c_iA_i)^2}{A_i}.}
  \label{eq:quotient-energy}
\end{equation}
Equation~\eqref{eq:quotient-energy} is the most convenient scalar expression for local automatic differentiation.  It also makes clear that a valid implementation requires $A_i>0$.

\section{Exact gradient and Hessian}

\subsection{Primitive derivative notation}

For each vertex energy $E_i$, define derivatives with respect to the relevant Cartesian coordinate vector, either global or local:
\begin{equation}
  \aV_i:=\grad A_i,
  \qquad
  \mV_i:=\grad M_i,
  \qquad
  \Q_i:=\Hess A_i,
  \qquad
  \Mmat_i:=\Hess M_i.
  \label{eq:primitive-derivatives}
\end{equation}
Define also
\begin{equation}
  \boxed{
  \gV_i:=\mV_i-H_i\aV_i.}
  \label{eq:g-def}
\end{equation}
Then
\begin{equation}
  \boxed{
  \grad H_i=\frac{\gV_i}{A_i}.}
  \label{eq:grad-H}
\end{equation}
The vector $\gV_i$ is the numerator of the curvature gradient and has dimensions of zero in length units: $M_i$ has length, so its gradient with respect to position is dimensionless; $H_i\aV_i$ is also dimensionless.

\subsection{Exact energy gradient}

Differentiating $E_i=\kappa_iA_i(H_i-c_i)^2$ while holding $\kappa_i$ and $c_i$ fixed gives
\begin{align}
  \grad E_i
  &=\kappa_i\left[
       2A_i(H_i-c_i)\grad H_i
       +(H_i-c_i)^2\aV_i
     \right]\\
  &=\kappa_i\left[
       2(H_i-c_i)(\mV_i-H_i\aV_i)
       +(H_i-c_i)^2\aV_i
     \right].
\end{align}
Therefore
\begin{equation}
  \boxed{
  \grad E_i
  =
  \kappa_i\left[
    2(H_i-c_i)\mV_i
    -(H_i^2-c_i^2)\aV_i
  \right].}
  \label{eq:local-gradient}
\end{equation}
Summing these local gradients and taking the negative reproduces the Mem3DG bending force.  Expanding $\mV_i$ into edge-length and dihedral-angle derivatives and expanding $\aV_i$ into triangle-area derivatives gives the halfedge-grouped expression in Appendix~C of the Mem3DG paper.

\subsection{Exact local Cartesian Hessian}

Differentiate Equation~\eqref{eq:local-gradient}.  Since
\begin{equation}
  \grad(H_i-c_i)=\frac{\gV_i}{A_i},
  \qquad
  \grad(H_i^2-c_i^2)=\frac{2H_i\gV_i}{A_i},
\end{equation}
the outer-product terms simplify:
\begin{align}
  &\frac{2}{A_i}\mV_i\gV_i^\T
  -\frac{2H_i}{A_i}\aV_i\gV_i^\T
  =\frac{2}{A_i}(\mV_i-H_i\aV_i)\gV_i^\T
  =\frac{2}{A_i}\gV_i\gV_i^\T.
\end{align}
The exact local Hessian is therefore
\begin{equation}
  \boxed{
  \Kmat_i
  :=\Hess E_i
  =
  \kappa_i\left[
    \frac{2}{A_i}\gV_i\gV_i^\T
    +2(H_i-c_i)\Mmat_i
    -(H_i^2-c_i^2)\Q_i
  \right].}
  \label{eq:master-hessian}
\end{equation}
The global bending Hessian is
\begin{equation}
  \boxed{
  \Kmat_b=\sum_i\Kmat_i,}
  \label{eq:global-hessian}
\end{equation}
with local-to-global scatter understood as in Equation~\eqref{eq:sparse-assembly}.

Equation~\eqref{eq:master-hessian} is the main implementation result.

\subsection{Interpretation of every term}

The three terms in Equation~\eqref{eq:master-hessian} have distinct meanings:
\begin{enumerate}
  \item \textbf{Curvature-gradient outer product}
  \begin{equation}
    \frac{2\kappa_i}{A_i}\gV_i\gV_i^\T
    =2\kappa_i A_i(\grad H_i)(\grad H_i)^\T.
  \end{equation}
  This term is positive semidefinite.  It is present even when $H_i=c_i$ and supplies the leading quadratic stiffness against changes of curvature.

  \item \textbf{Integrated-mean-curvature Hessian}
  \begin{equation}
    2\kappa_i(H_i-c_i)\Mmat_i.
  \end{equation}
  This term contains the second derivatives of edge lengths and signed dihedral angles.  It vanishes at a vertex where the discrete curvature exactly equals the spontaneous curvature.

  \item \textbf{Dual-area Hessian}
  \begin{equation}
    -\kappa_i(H_i^2-c_i^2)\Q_i.
  \end{equation}
  This term accounts for the shape dependence of the integration weight $A_i$ and for the area appearing in the denominator of $H_i=M_i/A_i$.
\end{enumerate}

The exact Hessian is generally indefinite away from a stable local minimum.  This is expected for a nonlinear geometric energy and is not, by itself, an implementation error.

\subsection{Equivalent formula using the Hessian of pointwise curvature}

For debugging, it is useful to record the explicit Hessian of $H_i=M_i/A_i$:
\begin{equation}
  \boxed{
  \Hess H_i
  =
  \frac{\Mmat_i}{A_i}
  -\frac{\mV_i\aV_i^\T+\aV_i\mV_i^\T}{A_i^2}
  -\frac{H_i}{A_i}\Q_i
  +\frac{2H_i}{A_i^2}\aV_i\aV_i^\T.}
  \label{eq:hess-H}
\end{equation}
A direct chain-rule expression is
\begin{align}
  \Hess E_i
  =\kappa_i\bigl[
    &(H_i-c_i)^2\Q_i
    +2(H_i-c_i)
      \bigl(\aV_i(\grad H_i)^\T+(\grad H_i)\aV_i^\T\bigr)\\
    &+2A_i(\grad H_i)(\grad H_i)^\T
    +2A_i(H_i-c_i)\Hess H_i
  \bigr].
  \label{eq:unsimplified-hessian}
\end{align}
Substituting Equation~\eqref{eq:hess-H} into Equation~\eqref{eq:unsimplified-hessian} simplifies exactly to Equation~\eqref{eq:master-hessian}.  The simplified formula is preferable in code because it has fewer terms and exposes symmetry.

\section{Exact Hessian--vector product}

For a global displacement direction $\vct\in\R^{3N}$, define
\begin{equation}
  \delta A_i=\aV_i^\T\vct,
  \qquad
  \delta M_i=\mV_i^\T\vct,
  \qquad
  \delta H_i=\frac{\delta M_i-H_i\delta A_i}{A_i}.
  \label{eq:directional-primitives}
\end{equation}
Since $\gV_i^\T\vct=A_i\delta H_i$, Equation~\eqref{eq:master-hessian} gives
\begin{equation}
  \boxed{
  \Kmat_i\vct
  =\kappa_i\left[
      2\delta H_i\gV_i
      +2(H_i-c_i)(\Mmat_i\vct)
      -(H_i^2-c_i^2)(\Q_i\vct)
    \right].}
  \label{eq:hvp}
\end{equation}
This expression is ideal for Newton--Krylov methods.  It requires only directional second derivatives
\begin{equation}
  \Q_i\vct=D(\grad A_i)[\vct],
  \qquad
  \Mmat_i\vct=D(\grad M_i)[\vct],
\end{equation}
not the full dense local matrices.

\section{Primitive derivative formulas}

\subsection{A reusable normalized-vector identity}

Many geometric quantities are normalized vectors.  Let
\begin{equation}
  \widehat{\mathbf{u}}=\frac{\mathbf{u}}{\rho},
  \qquad
  \rho=\|\mathbf{u}\|,
  \qquad
  \Pmat_{\mathbf{u}}=\Id-\widehat{\mathbf{u}}\widehat{\mathbf{u}}^\T.
\end{equation}
For perturbation directions $\mathbf{p}$ and $\mathbf{q}$ in the underlying coordinates,
\begin{equation}
  \boxed{
  D\widehat{\mathbf{u}}[\mathbf{p}]
  =\frac{\Pmat_{\mathbf{u}}D\mathbf{u}[\mathbf{p}]}{\rho}.}
  \label{eq:normalized-first}
\end{equation}
The second derivative is
\begin{equation}
  \boxed{\begin{aligned}
  D^2\widehat{\mathbf{u}}[\mathbf{p},\mathbf{q}]
  ={}&\frac{\Pmat_{\mathbf{u}}D^2\mathbf{u}[\mathbf{p},\mathbf{q}]}{\rho}
  \\
  &-\frac{1}{\rho^2}\Bigl[
    (\widehat{\mathbf{u}}\cdot D\mathbf{u}[\mathbf{p}])
      \Pmat_{\mathbf{u}}D\mathbf{u}[\mathbf{q}]
  \\
  &\hspace{19mm}+
    (\widehat{\mathbf{u}}\cdot D\mathbf{u}[\mathbf{q}])
      \Pmat_{\mathbf{u}}D\mathbf{u}[\mathbf{p}]
  \\
  &\hspace{19mm}+
    \widehat{\mathbf{u}}
      \bigl(D\mathbf{u}[\mathbf{p}]^\T
      \Pmat_{\mathbf{u}}D\mathbf{u}[\mathbf{q}]\bigr)
  \Bigr].
  \end{aligned}}
  \label{eq:normalized-second}
\end{equation}
This identity is used for edge tangents and face normals.

\subsection{Edge length}

For edge $e=(p,q)$, define
\begin{equation}
  \mathbf{d}=\mathbf{x}_q-\mathbf{x}_p,
  \qquad
  \ell=\|\mathbf{d}\|,
  \qquad
  \mathbf{t}=\frac{\mathbf{d}}{\ell},
  \qquad
  \Pmat_t=\Id-\mathbf{t}\mathbf{t}^\T.
\end{equation}
In local ordering $(\mathbf{x}_p,\mathbf{x}_q)$,
\begin{equation}
  \boxed{
  \grad\ell=
  \begin{bmatrix}
    -\mathbf{t}\\
     \mathbf{t}
  \end{bmatrix},}
\end{equation}
and
\begin{equation}
  \boxed{
  \Hess\ell
  =\frac{1}{\ell}
  \begin{bmatrix}
    \Pmat_t&-\Pmat_t\\
    -\Pmat_t&\Pmat_t
  \end{bmatrix}.}
  \label{eq:edge-length-hessian}
\end{equation}
The Hessian has translational null modes and is positive semidefinite.

The edge tangent derivatives follow from Equations~\eqref{eq:normalized-first}--\eqref{eq:normalized-second} with $\mathbf{u}=\mathbf{d}$.  Because $\mathbf{d}$ is affine in the vertex coordinates, $D^2\mathbf{d}=0$.

\subsection{Triangle area}

For an oriented face $f=(p,q,r)$, set
\begin{equation}
  \mathbf{u}=\mathbf{x}_q-\mathbf{x}_p,
  \qquad
  \mathbf{v}=\mathbf{x}_r-\mathbf{x}_p,
  \qquad
  \mathbf{c}=\mathbf{u}\times\mathbf{v},
\end{equation}
\begin{equation}
  s=\|\mathbf{c}\|,
  \qquad
  A_f=\frac{s}{2},
  \qquad
  \mathbf{n}=\frac{\mathbf{c}}{s},
  \qquad
  \Pmat_n=\Id-\mathbf{n}\mathbf{n}^\T.
\end{equation}
For a coordinate perturbation $\mathbf{p}$, write
\begin{equation}
  \delta\mathbf{u}_{\mathbf{p}}
  =\delta\mathbf{x}_{q,\mathbf{p}}-\delta\mathbf{x}_{p,\mathbf{p}},
  \qquad
  \delta\mathbf{v}_{\mathbf{p}}
  =\delta\mathbf{x}_{r,\mathbf{p}}-\delta\mathbf{x}_{p,\mathbf{p}}.
\end{equation}
Then
\begin{equation}
  D\mathbf{c}[\mathbf{p}]
  =\delta\mathbf{u}_{\mathbf{p}}\times\mathbf{v}
   +\mathbf{u}\times\delta\mathbf{v}_{\mathbf{p}},
  \label{eq:cross-first}
\end{equation}
and
\begin{equation}
  D^2\mathbf{c}[\mathbf{p},\mathbf{q}]
  =\delta\mathbf{u}_{\mathbf{p}}\times\delta\mathbf{v}_{\mathbf{q}}
   +\delta\mathbf{u}_{\mathbf{q}}\times\delta\mathbf{v}_{\mathbf{p}}.
  \label{eq:cross-second}
\end{equation}
The first variation is
\begin{equation}
  DA_f[\mathbf{p}]
  =\frac{1}{2}\mathbf{n}\cdot D\mathbf{c}[\mathbf{p}].
\end{equation}
The standard vertex gradients are
\begin{align}
  \grad_{\mathbf{x}_p}A_f
  &=\frac{1}{2}(\mathbf{x}_q-\mathbf{x}_r)\times\mathbf{n},\\
  \grad_{\mathbf{x}_q}A_f
  &=\frac{1}{2}(\mathbf{x}_r-\mathbf{x}_p)\times\mathbf{n},\\
  \grad_{\mathbf{x}_r}A_f
  &=\frac{1}{2}(\mathbf{x}_p-\mathbf{x}_q)\times\mathbf{n}.
\end{align}
The exact area Hessian is most compactly expressed as a bilinear form:
\begin{equation}
  \boxed{
  D^2A_f[\mathbf{p},\mathbf{q}]
  =\frac{1}{2}\left[
    \frac{D\mathbf{c}[\mathbf{p}]^\T\Pmat_nD\mathbf{c}[\mathbf{q}]}{s}
    +\mathbf{n}\cdot D^2\mathbf{c}[\mathbf{p},\mathbf{q}]
  \right].}
  \label{eq:triangle-area-hessian}
\end{equation}
Equation~\eqref{eq:triangle-area-hessian} defines the $9\times9$ local face Hessian exactly.  It can be converted to a matrix by applying it to coordinate basis directions, or implemented directly through Jacobians of $\mathbf{c}$.

An equivalent matrix formula is
\begin{equation}
  \Hess A_f
  =\frac{1}{2}\left[
    \frac{\Jmat_c^\T\Pmat_n\Jmat_c}{s}
    +\sum_{\alpha=1}^{3}n_\alpha\Hess c_\alpha
  \right],
  \label{eq:area-hessian-matrix}
\end{equation}
where $\Jmat_c$ is the $3\times9$ Jacobian of the cross-product vector and $\Hess c_\alpha$ is the $9\times9$ Hessian of its $\alpha$th component.

The face-normal derivatives follow from Equations~\eqref{eq:normalized-first}--\eqref{eq:normalized-second} with $\mathbf{u}=\mathbf{c}$, using Equations~\eqref{eq:cross-first}--\eqref{eq:cross-second}.

\subsection{Dual-area gradient and Hessian}

By linearity of Equation~\eqref{eq:dual-area},
\begin{equation}
  \boxed{
  \aV_i=\frac{1}{3}\sum_{f\ni i}\grad A_f,
  \qquad
  \Q_i=\frac{1}{3}\sum_{f\ni i}\Hess A_f.}
  \label{eq:dual-area-derivatives}
\end{equation}
Every face derivative is scattered into the local one-ring coordinates of vertex $i$.

\subsection{Signed dihedral-angle gradient and Hessian}

Use the quantities in Equation~\eqref{eq:signed-dihedral} and define
\begin{equation}
  r^2=\sigma^2+\chi^2.
\end{equation}
For exact unit normals, $r^2=1$ analytically.  Keeping $r^2$ in the formulas improves robustness and makes the chain rule explicit.

Let
\begin{equation}
  \mathbf{g}_{\sigma}=\grad\sigma,
  \qquad
  \mathbf{g}_{\chi}=\grad\chi,
  \qquad
  \mathbf{H}_{\sigma}=\Hess\sigma,
  \qquad
  \mathbf{H}_{\chi}=\Hess\chi.
\end{equation}
Then
\begin{equation}
  \boxed{
  \grad\phi
  =\frac{\chi\mathbf{g}_{\sigma}-\sigma\mathbf{g}_{\chi}}{r^2}.}
  \label{eq:dihedral-gradient}
\end{equation}
The exact dihedral Hessian is
\begin{equation}
  \boxed{\begin{aligned}
  \Hess\phi
  ={}&\frac{\chi}{r^2}\mathbf{H}_{\sigma}
     -\frac{\sigma}{r^2}\mathbf{H}_{\chi}
  \\
  &-\frac{2\sigma\chi}{r^4}
      \mathbf{g}_{\sigma}\mathbf{g}_{\sigma}^\T
  +\frac{\sigma^2-\chi^2}{r^4}
      \left(
        \mathbf{g}_{\sigma}\mathbf{g}_{\chi}^\T
        +\mathbf{g}_{\chi}\mathbf{g}_{\sigma}^\T
      \right)
  \\
  &+\frac{2\sigma\chi}{r^4}
      \mathbf{g}_{\chi}\mathbf{g}_{\chi}^\T.
  \end{aligned}}
  \label{eq:dihedral-hessian}
\end{equation}
This formula is the Hessian chain rule for $\atanTwo(\sigma,\chi)$.

The first directional derivatives needed to construct $\mathbf{g}_{\sigma}$ and $\mathbf{g}_{\chi}$ are
\begin{align}
  D\sigma[\mathbf{p}]
  ={}&D\mathbf{t}[\mathbf{p}]\cdot(\mathbf{n}_+\times\mathbf{n}_-)
  \\
  &+\mathbf{t}\cdot\left(
      D\mathbf{n}_+[\mathbf{p}]\times\mathbf{n}_-
      +\mathbf{n}_+\times D\mathbf{n}_-[\mathbf{p}]
    \right),
  \label{eq:sigma-first}\\
  D\chi[\mathbf{p}]
  ={}&D\mathbf{n}_+[\mathbf{p}]\cdot\mathbf{n}_-
      +\mathbf{n}_+\cdot D\mathbf{n}_-[\mathbf{p}].
  \label{eq:chi-first}
\end{align}
The second directional derivatives are
\begin{align}
  D^2\sigma[\mathbf{p},\mathbf{q}]
  ={}&D^2\mathbf{t}[\mathbf{p},\mathbf{q}]
       \cdot(\mathbf{n}_+\times\mathbf{n}_-)
  \\
  &+D\mathbf{t}[\mathbf{p}]\cdot\left(
       D\mathbf{n}_+[\mathbf{q}]\times\mathbf{n}_-
       +\mathbf{n}_+\times D\mathbf{n}_-[\mathbf{q}]
     \right)
  \\
  &+D\mathbf{t}[\mathbf{q}]\cdot\left(
       D\mathbf{n}_+[\mathbf{p}]\times\mathbf{n}_-
       +\mathbf{n}_+\times D\mathbf{n}_-[\mathbf{p}]
     \right)
  \\
  &+\mathbf{t}\cdot\Bigl(
       D^2\mathbf{n}_+[\mathbf{p},\mathbf{q}]\times\mathbf{n}_-
       +D\mathbf{n}_+[\mathbf{p}]\times D\mathbf{n}_-[\mathbf{q}]
  \\
  &\hspace{25mm}
       +D\mathbf{n}_+[\mathbf{q}]\times D\mathbf{n}_-[\mathbf{p}]
       +\mathbf{n}_+\times D^2\mathbf{n}_-[\mathbf{p},\mathbf{q}]
     \Bigr),
  \label{eq:sigma-second}\\
  D^2\chi[\mathbf{p},\mathbf{q}]
  ={}&D^2\mathbf{n}_+[\mathbf{p},\mathbf{q}]\cdot\mathbf{n}_-
  \\
  &+D\mathbf{n}_+[\mathbf{p}]\cdot D\mathbf{n}_-[\mathbf{q}]
   +D\mathbf{n}_+[\mathbf{q}]\cdot D\mathbf{n}_-[\mathbf{p}]
  \\
  &+\mathbf{n}_+\cdot D^2\mathbf{n}_-[\mathbf{p},\mathbf{q}].
  \label{eq:chi-second}
\end{align}
The required derivatives of $\mathbf{t}$ and $\mathbf{n}_{\pm}$ are all instances of Equations~\eqref{eq:normalized-first}--\eqref{eq:normalized-second}.

For production code, Equation~\eqref{eq:dihedral-hessian} is often best implemented by second-order automatic differentiation of the small four-vertex edge-diamond function.  It is also reasonable to use a specialized closed-form hinge-angle Hessian, such as the formulas of Tamstorf and Grinspun, provided its orientation and angle convention exactly match the Mem3DG convention.

\subsection{Integrated mean-curvature gradient and Hessian}

From Equation~\eqref{eq:integrated-mean-curvature},
\begin{equation}
  M_i=\frac{1}{4}\sum_{e\ni i}\ell_e\phi_e.
\end{equation}
Therefore
\begin{equation}
  \boxed{
  \mV_i
  =\frac{1}{4}\sum_{e\ni i}
    \left(
      \phi_e\grad\ell_e
      +\ell_e\grad\phi_e
    \right).}
  \label{eq:M-gradient}
\end{equation}
The exact Hessian is
\begin{equation}
  \boxed{\begin{aligned}
  \Mmat_i
  =\frac{1}{4}\sum_{e\ni i}\Bigl[
     &\phi_e\Hess\ell_e
      +\ell_e\Hess\phi_e
  \\
     &+(\grad\ell_e)(\grad\phi_e)^\T
      +(\grad\phi_e)(\grad\ell_e)^\T
  \Bigr].
  \end{aligned}}
  \label{eq:M-hessian}
\end{equation}
The last two terms are required by the product rule and are essential for symmetry.  Omitting either one produces an incorrect Jacobian.

Each edge term is naturally computed on its four-vertex diamond and scattered into the one-ring coordinate vector for $M_i$.  Boundary edges require a boundary-specific curvature convention; a common closed-surface implementation simply excludes boundary dihedral terms.

\section{Recommended implementation strategies}

\subsection{Strategy A: local scalar energy with second-order automatic differentiation}

This is the recommended first implementation because it is compact and difficult to get algebraically wrong.

For each center vertex $i$:
\begin{enumerate}
  \item Gather the positions of the closed one-ring $S_i$ into a local vector $\x_i^{\mathrm{loc}}$.
  \item Recompute all incident triangle areas and all incident signed dihedral angles using only the local coordinates.
  \item Compute
  \begin{equation}
    A_i=\frac{1}{3}\sum_{f\ni i}A_f,
    \qquad
    M_i=\frac{1}{4}\sum_{e\ni i}\ell_e\phi_e.
  \end{equation}
  \item Evaluate the scalar local energy
  \begin{equation}
    E_i=\kappa_i\frac{(M_i-c_iA_i)^2}{A_i}.
  \end{equation}
  \item Ask the AD library for the local gradient and local Hessian.
  \item Scatter-add them into the global gradient and sparse Hessian.
\end{enumerate}

A C++-like sketch is:
\begin{lstlisting}[language=C++]
for (Vertex i : mesh.vertices()) {
    LocalStencil S = buildClosedOneRingStencil(i);
    ADVector x = gatherADPositions(S);

    ADScalar A = 0.0;
    for (Face f : incidentFaces(i)) {
        A += triangleAreaAD(x, f, S) / 3.0;
    }

    ADScalar M = 0.0;
    for (Edge e : incidentEdges(i)) {
        ADScalar l   = edgeLengthAD(x, e, S);
        ADScalar phi = signedDihedralAD(x, e, S);
        M += 0.25 * l * phi;
    }

    // For the current Mem3DG source use kappa = 2 * Kb[i].
    ADScalar E = kappa[i] * square(M - c[i] * A) / A;

    Vector gloc = gradient(E, x);
    Matrix Kloc = hessian(E, x);

    scatterAdd(globalGradient, gloc, S);
    scatterAdd(globalHessian,  Kloc, S);
}
\end{lstlisting}

Advantages:
\begin{itemize}
  \item It differentiates exactly the scalar energy actually implemented.
  \item Symmetry follows automatically, up to floating-point and AD tolerances.
  \item It avoids differentiating the already-grouped halfedge force expression.
  \item It provides a high-quality reference against which a hand-optimized analytic implementation can be tested.
\end{itemize}

Disadvantages:
\begin{itemize}
  \item Second-order AD can be expensive if the local tape is rebuilt repeatedly.
  \item Dynamic one-ring sizes complicate template-based fixed-size AD.
  \item A generic AD package may allocate heavily unless local buffers are reused.
\end{itemize}

\subsection{Strategy B: analytic primitive Hessians plus the master formula}

For each vertex $i$:
\begin{enumerate}
  \item Accumulate $A_i$, $\aV_i$, and $\Q_i$ from incident face areas.
  \item Accumulate $M_i$, $\mV_i$, and $\Mmat_i$ from incident edge products $\ell_e\phi_e$.
  \item Compute $H_i=M_i/A_i$ and $\gV_i=\mV_i-H_i\aV_i$.
  \item Evaluate Equation~\eqref{eq:local-gradient} and Equation~\eqref{eq:master-hessian}.
  \item Scatter-add to global arrays.
\end{enumerate}

A C++-like sketch is:
\begin{lstlisting}[language=C++]
for (Vertex i : mesh.vertices()) {
    Stencil S = closedOneRing(i);
    double A = 0.0;
    Vector a = zeros(3 * S.size());
    Matrix Q = zeros(3 * S.size(), 3 * S.size());

    for (Face f : incidentFaces(i)) {
        AreaDerivatives df = triangleAreaDerivatives(f, S);
        A += df.value / 3.0;
        a += df.gradient / 3.0;
        Q += df.hessian / 3.0;
    }

    double M = 0.0;
    Vector m = zerosLike(a);
    Matrix R = zerosLike(Q);

    for (Edge e : incidentEdges(i)) {
        EdgeLengthDerivatives dl = edgeLengthDerivatives(e, S);
        DihedralDerivatives  dp = signedDihedralDerivatives(e, S);

        M += 0.25 * dl.value * dp.value;
        m += 0.25 * (dp.value * dl.gradient
                   + dl.value * dp.gradient);
        R += 0.25 * (dp.value * dl.hessian
                   + dl.value * dp.hessian
                   + outer(dl.gradient, dp.gradient)
                   + outer(dp.gradient, dl.gradient));
    }

    double H = M / A;
    Vector g = m - H * a;
    double d = H - c[i];

    Vector gradE = kappa[i] * (2.0 * d * m
                             - (H * H - c[i] * c[i]) * a);

    Matrix hessE = kappa[i] * (
          (2.0 / A) * outer(g, g)
        + 2.0 * d * R
        - (H * H - c[i] * c[i]) * Q);

    scatterAdd(globalGradient, gradE, S);
    scatterAdd(globalHessian, hessE, S);
}
\end{lstlisting}

This organization is more maintainable than differentiating the grouped Mem3DG force term by term.  The latter requires derivatives of the mean-curvature, Gaussian-curvature, and Schlaefli vectors and tends to obscure exact Hessian symmetry.

\subsection{Strategy C: matrix-free directional differentiation}

For a Hessian--vector product, compute the directional derivatives of every local primitive under the gathered local displacement direction.  Use Equation~\eqref{eq:hvp}.  A forward-over-reverse AD mode on the local scalar energy is often efficient for this purpose.

\section{Sparse matrix assembly details}

\subsection{Local indexing}

Create a deterministic ordered list
\begin{equation}
  S_i=(s_0,s_1,\ldots,s_{m_i-1})
\end{equation}
with $s_0=i$ if convenient.  Maintain a map from global vertex index to local stencil index.  Every face or edge-diamond derivative is first expressed in its natural local ordering and then scattered into the $3m_i$ local vector or matrix.

For a local vertex index $a$, its three coordinate columns are
\begin{equation}
  3a,
  \quad 3a+1,
  \quad 3a+2.
\end{equation}
Use the same convention globally.

\subsection{Triplet assembly}

For an assembled sparse Hessian, append every nonzero $3\times3$ block of each $\Kmat_i^{\mathrm{loc}}$ as sparse triplets.  Let the sparse-matrix constructor sum duplicate entries.  Do not overwrite duplicate blocks, because many vertex energies contribute to the same global pair.

After assembly, numerical symmetrization
\begin{equation}
  \Kmat_b\leftarrow\frac{1}{2}(\Kmat_b+\Kmat_b^\T)
\end{equation}
may remove roundoff asymmetry, but it must not be used to hide a derivation error.  The unsymmetrized relative defect should already be near machine precision for an analytic or AD implementation.

\subsection{Constraint handling}

Apply fixed-vertex or fixed-coordinate constraints after assembling the unconstrained Hessian.  Common choices are:
\begin{itemize}
  \item eliminate constrained degrees of freedom;
  \item replace constrained rows and columns by identity; or
  \item solve in a projected free subspace.
\end{itemize}
The chosen operation must be applied consistently to the residual or force vector.

\section{Verification and regression tests}

A Hessian implementation should not be accepted based only on visual simulation behavior.  Use the tests below on random nondegenerate meshes, a tetrahedron, an icosphere, and meshes with heterogeneous $\kappa_i$ and $c_i$.

\subsection{Energy--gradient directional test}

Choose a random normalized direction $\vct$.  Verify
\begin{equation}
  \frac{E(\x+\varepsilon\vct)-E(\x-\varepsilon\vct)}{2\varepsilon}
  \approx
  \grad E(\x)^\T\vct.
  \label{eq:grad-test}
\end{equation}
For a range of $\varepsilon$, the error should initially decrease like $O(\varepsilon^2)$ and then increase due to roundoff.

\subsection{Gradient--Hessian directional test}

Verify
\begin{equation}
  \frac{\grad E(\x+\varepsilon\vct)-\grad E(\x-\varepsilon\vct)}{2\varepsilon}
  \approx
  \Kmat_b(\x)\vct.
  \label{eq:hessian-test}
\end{equation}
This is the most important implementation test.  It checks the complete Hessian without requiring a finite-difference approximation of every matrix entry.

\subsection{Force-Jacobian test}

If the code stores physical force $\mathbf{f}=-\grad E$, verify
\begin{equation}
  \frac{\mathbf{f}(\x+\varepsilon\vct)-\mathbf{f}(\x-\varepsilon\vct)}{2\varepsilon}
  \approx
  -\Kmat_b\vct.
\end{equation}
A frequent error is a globally correct Hessian with the wrong sign in the Newton system.

\subsection{Symmetry test}

Compute
\begin{equation}
  \eta_{\mathrm{sym}}
  =\frac{\|\Kmat_b-\Kmat_b^\T\|_F}
         {\max(1,\|\Kmat_b\|_F)}.
\end{equation}
For a double-precision analytic or AD Hessian, $\eta_{\mathrm{sym}}$ should typically be near $10^{-12}$ to $10^{-14}$, depending on mesh conditioning and assembly order.

\subsection{Translation invariance}

The energy is invariant under translation.  For each Cartesian basis vector $\mathbf{e}_\alpha$, define a global translation direction
\begin{equation}
  \vct_{\mathrm{tr},\alpha}
  =(\mathbf{e}_\alpha,\ldots,\mathbf{e}_\alpha).
\end{equation}
Then
\begin{equation}
  \grad E^\T\vct_{\mathrm{tr},\alpha}=0,
  \qquad
  \Kmat_b\vct_{\mathrm{tr},\alpha}=\zero.
\end{equation}
The second identity holds because the translation path is affine and the energy is exactly constant along it.

\subsection{Rotation check and its caveat}

The energy is invariant under rigid rotation, but the infinitesimal rotational direction $\delta\x=\bm{\omega}\times\x$ is not an affine exact rotation path at second order.  Away from equilibrium, one should not expect $\Kmat_b\delta\x=0$ by itself.  At a critical configuration where $\grad E=0$, the three rotational modes become Hessian null modes up to numerical error.

\subsection{Local formula cross-check}

For each vertex $i$, compare:
\begin{enumerate}
  \item the Hessian obtained by second-order AD of the scalar $E_i$; and
  \item the Hessian from Equation~\eqref{eq:master-hessian}, using AD only to obtain $\aV_i,\mV_i,\Q_i,\Mmat_i$.
\end{enumerate}
They should agree to roundoff.  This is a strong algebraic test of the master formula.

\subsection{Normalization cross-check against Mem3DG}

To match the current Mem3DG source, use $\kappa_i=2K_{b,i}$.  Verify that
\begin{equation}
  -\grad E_b
\end{equation}
matches the existing spontaneous-curvature force before enabling the Hessian in an implicit solver.  First test homogeneous $K_b$ and $H_0=0$, then heterogeneous fields.

\section{Degeneracy, robustness, and solver behavior}

\subsection{Geometric degeneracy}

The formulas contain divisions by
\begin{equation}
  \ell_e,
  \qquad
  \|\mathbf{c}_f\|=2A_f,
  \qquad
  A_i.
\end{equation}
A triangle approaching zero area causes normal, dihedral, gradient, and Hessian blow-up.  Do not silently clamp these values inside the differentiated energy unless the regularized energy is intentionally changed.  Better options are:
\begin{itemize}
  \item reject invalid Newton trial steps;
  \item use line search or trust-region step control;
  \item maintain mesh quality through remeshing outside the derivative evaluation; or
  \item define and consistently differentiate an explicit regularized energy.
\end{itemize}

\subsection{Indefinite Newton systems}

The exact Hessian may be indefinite.  Stable nonlinear solves commonly use one of:
\begin{itemize}
  \item damped Newton with backtracking line search;
  \item trust-region Newton;
  \item an LDL$^\T$ factorization with inertia control;
  \item a shifted matrix $\Kmat_b+\mu\Id$; or
  \item a positive-semidefinite approximation for selected terms.
\end{itemize}
Do not replace the exact Hessian by its absolute eigenvalue projection without documenting that the method is no longer exact Newton.

\subsection{Caching}

Any cached edge length, face normal, face area, dihedral angle, dual area, or integrated curvature must correspond to the same coordinate state at which the Hessian is evaluated.  A stale geometry cache can pass some energy tests while failing the Hessian directional test.

\section{Gaussian curvature and boundaries}

For a closed mesh with fixed topology and constant Gaussian modulus $\kappa_G$,
\begin{equation}
  \sum_i\int K_i
  =2\pi\chi(\mathcal{M})
\end{equation}
is constant by the discrete Gauss--Bonnet theorem.  Therefore
\begin{equation}
  \grad E_G=\zero,
  \qquad
  \Hess E_G=\zero.
\end{equation}

This conclusion does not apply when:
\begin{itemize}
  \item the mesh has a boundary;
  \item $\kappa_G$ varies spatially;
  \item topology changes; or
  \item the discretization assigns a non-topological Gaussian energy.
\end{itemize}
In those cases, the angle-defect and boundary geodesic-curvature terms require their own gradients and Hessians.

\section{Geometry-dependent material fields}

Equation~\eqref{eq:master-hessian} assumes $\kappa_i$ and $c_i$ are fixed during the shape derivative.  This is the correct interpretation when they are independent state variables held fixed in a mechanical Newton solve.

If $\kappa_i=\kappa_i(\x)$ or $c_i=c_i(\x)$, the energy is still
\begin{equation}
  E_i=\kappa_iA_i(H_i-c_i)^2,
\end{equation}
but the gradient gains
\begin{equation}
  A_i(H_i-c_i)^2\grad\kappa_i
  -2\kappa_iA_i(H_i-c_i)\grad c_i,
\end{equation}
and the Hessian gains all product-rule derivatives of these terms.  In a coupled mechanochemical system, it is usually cleaner to assemble the full block Jacobian in variables $(\x,\phi)$ rather than hiding chemical dependence inside a shape-only Hessian.

\section{Dimensional and scaling checks}

With position measured in length units $L$:
\begin{center}
\begin{tabular}{lll}
\toprule
Quantity & Dimension & Comment\\
\midrule
$A_i$ & $L^2$ & dual area\\
$M_i$ & $L$ & integrated mean curvature\\
$H_i,c_i$ & $L^{-1}$ & pointwise curvature\\
$\kappa_i$ & energy & bending modulus in Equation~\eqref{eq:paper-energy}\\
$E_i$ & energy & $\kappa_iA_i(H_i-c_i)^2$\\
$\grad E_i$ & energy/$L$ & force magnitude\\
$\Hess E_i$ & energy/$L^2$ & stiffness\\
\bottomrule
\end{tabular}
\end{center}
A uniform scaling $\x\mapsto s\x$ gives $A_i\mapsto s^2A_i$, $M_i\mapsto sM_i$, and $H_i\mapsto H_i/s$.  For $c_i=0$, the pure Willmore part $\sum_i\kappa_iA_iH_i^2$ is scale invariant.  This provides another useful regression check.

\section{Implementation checklist for Codex or a human developer}

\begin{enumerate}
  \item Confirm whether the code expects the energy Hessian $\Hess E$ or the force Jacobian $-\Hess E$.
  \item Confirm the energy normalization.  Current Mem3DG source: $\kappa_i=2K_{b,i}$.
  \item Treat \texttt{vertexMeanCurvatures} as $M_i$, not $H_i$.
  \item Use a single signed-dihedral convention everywhere.
  \item Build $A_i$ from one-third of incident face areas.
  \item Build $M_i$ from one-quarter of $\ell_e\phi_e$ over incident edges.
  \item Prefer local scalar second-order AD for the first correct implementation.
  \item If implementing analytically, include both mixed outer products in Equation~\eqref{eq:M-hessian}.
  \item Assemble duplicate sparse contributions by summation.
  \item Test energy--gradient consistency.
  \item Test gradient--Hessian consistency.
  \item Test symmetry and translation null modes.
  \item Compare $-\grad E$ against the existing Mem3DG bending force.
  \item Use line search or trust-region protection against triangle inversion and dihedral branch jumps.
\end{enumerate}

\section{Concise formula sheet}

For each vertex $i$:
\begin{align}
  A_i&=\frac13\sum_{f\ni i}A_f,\\
  M_i&=\frac14\sum_{e\ni i}\ell_e\phi_e,\\
  H_i&=\frac{M_i}{A_i},\\
  E_i&=\kappa_iA_i(H_i-c_i)^2,\\
  \aV_i&=\grad A_i,
  &\mV_i&=\grad M_i,\\
  \Q_i&=\Hess A_i,
  &\Mmat_i&=\Hess M_i,\\
  \gV_i&=\mV_i-H_i\aV_i.
\end{align}
Then
\begin{equation}
  \boxed{
  \grad E_i
  =\kappa_i\left[
    2(H_i-c_i)\mV_i
    -(H_i^2-c_i^2)\aV_i
  \right],}
\end{equation}
\begin{equation}
  \boxed{
  \Hess E_i
  =\kappa_i\left[
    \frac{2}{A_i}\gV_i\gV_i^\T
    +2(H_i-c_i)\Mmat_i
    -(H_i^2-c_i^2)\Q_i
  \right],}
\end{equation}
and
\begin{equation}
  \boxed{
  \frac{\partial\mathbf{f}_b}{\partial\x}
  =-\sum_i\Hess E_i.}
\end{equation}

The primitive integrated-curvature derivatives are
\begin{equation}
  \boxed{
  \mV_i
  =\frac14\sum_{e\ni i}
  \left(\phi_e\grad\ell_e+\ell_e\grad\phi_e\right),}
\end{equation}
\begin{equation}
  \boxed{
  \Mmat_i
  =\frac14\sum_{e\ni i}
  \left[
    \phi_e\Hess\ell_e
    +\ell_e\Hess\phi_e
    +(\grad\ell_e)(\grad\phi_e)^\T
    +(\grad\phi_e)(\grad\ell_e)^\T
  \right].}
\end{equation}

\section{References}

\begin{thebibliography}{9}

\bibitem{mem3dg-paper}
C. Zhu, C. T. Lee, and P. Rangamani,
``Mem3DG: Modeling membrane mechanochemical dynamics in 3D using discrete differential geometry,''
\emph{Biophysical Reports}, vol. 2, 100062, 2022.
\url{https://doi.org/10.1016/j.bpr.2022.100062}

\bibitem{mem3dg-repo}
RangamaniLabUCSD,
``Mem3DG: 3D computational model for lipid membranes,''
source repository.
\url{https://github.com/RangamaniLabUCSD/Mem3DG}

\bibitem{tamstorf-grinspun}
R. Tamstorf and E. Grinspun,
``Discrete bending forces and their Jacobians,''
\emph{Graphical Models}, vol. 75, no. 6, pp. 362--370, 2013.

\bibitem{crane-ddg}
K. Crane, F. de Goes, M. Desbrun, and P. Schroeder,
``Digital Geometry Processing with Discrete Exterior Calculus,''
SIGGRAPH Course Notes, 2013.

\end{thebibliography}

\end{document}
